\documentclass[a4paper,14pt]{extarticle}
\usepackage{suai-report}

\suaisetup{
  department   = 41,
  teacher-post = Старший преподаватель,
  teacher      = В. В. Боженко,
  type         = Отчет о лабораторной работе,
  number       = 5,
  title        = Классификация,
  course       = Введение в анализ данных,
  group        = 4414,
  student      = Г. В. Моисеев,
  date         = today,
}

\begin{document}

\maketitle
\suaitoc

\suaiintro

Цель работы — изучить методы классификации и сравнить качество предсказания целевого класса на практическом наборе данных.

\suaitasks
загрузить и проверить набор данных, выделить целевой признак и выполнить предварительную обработку
разделить данные на тренировочную и тестовую выборки, закодировать категориальные признаки и стандартизовать числовые
обучить классификаторы методом k-ближайших соседей, деревом решений, логистической регрессией и случайным лесом
вычислить метрики качества, построить матрицы неточностей и сравнить ROC-кривые моделей
сделать выводы о качестве классификации с учётом распределения классов

\section{Вариант задания}

Выполнен вариант 1 — набор данных \texttt{medical\_diagnosis.csv}. Каждая строка описывает пациента. Целевой признак \texttt{Disease}: 1 — пациент болеет, 0 — не болеет. Признак не включается в матрицу входных данных. Числовые признаки — возраст, индекс массы тела, артериальное давление, холестерин, глюкоза, выраженность симптомов, C-реактивный белок и число лейкоцитов. Категориальные признаки — пол, статус курения, наследственность, уровень активности и результат ЭКГ.

\suaitable[fields]{Признаки набора данных}
Признак | Содержание
age | возраст пациента, лет
sex | пол: Male или Female
BMI | индекс массы тела, кг/м²
systolic\_bp | систолическое давление, мм рт. ст.
cholesterol\_mg\_dl | холестерин, мг/дл
fasting\_glucose\_mg\_dl | глюкоза натощак, мг/дл
smoking\_status | статус курения: Never, Former или Current
family\_history | наследственность: Yes или No
activity\_level | физическая активность: Low, Moderate или High
symptom\_score | выраженность симптомов, от 0 до 10
CRP\_mg\_L | C-реактивный белок, мг/л
WBC\_10e9\_L | лейкоциты, $\times 10^9$/л
ECG\_result | ЭКГ: Normal, ST-T abnormality или LVH
Disease | целевой класс: 1 — болеет, 0 — не болеет

Работа выполнена в Jupyter-ноутбуке \texttt{ЛР5\_Моисеев\_1.ipynb}. В нём приведены пояснения, код, результаты обучения моделей и построения графиков.

\section{Исследование и обработка данных}

Исходная таблица содержит 3031 строку и 14 столбцов, пропусков нет. До обработки классу 0 принадлежат 2982 строки, классу 1 — 49. Найден один полный дубликат.

В конце файла подряд расположена группа из 31 строки, в которой холестерин превышает 300~мг/дл, а C-реактивный белок — 20~мг/л. В остальной части файла такого сочетания показателей нет; 26 значений холестерина в этой группе дробные, тогда как в остальных строках они целые. Группа содержит полный дубликат. Сочетание расположения в файле и отличающегося масштаба двух показателей указывает на загрязнённые записи. Дубликат и 30 уникальных строк этой группы исключены из анализа; остальные наблюдения не удалялись по правилу выбросов. До и после очистки данные показаны на~\figref{data_audit}.

\suaiimg[0.9\textwidth]{data_audit}{Распределение классов и подозрительная группа записей}

После обработки осталось 3000 строк: 2952 пациента класса 0 и 48 пациентов класса 1. Доля положительного класса равна 1,6~\%, то есть классы сильно несбалансированы. Для построения матрицы диаграмм рассеяния отображены все 48 пациентов класса 1 и случайная выборка из 600 пациентов класса 0; для обучения использованы все оставшиеся строки. Матрица построена по восьми числовым признакам и показана на~\figref{pairplot}.

\suaiimg[0.9\textwidth]{pairplot}{Матрица диаграмм рассеяния числовых признаков}

\section{Подготовка выборок и моделей}

Данные разделены на тренировочную и тестовую выборки в отношении 80:20 с параметрами \texttt{random\_state = 42} и \texttt{stratify = y}. В тренировочную выборку вошли 2400 строк (2362 класса 0 и 38 класса 1), в тестовую — 600 строк (590 класса 0 и 10 класса 1). Стратификация сохраняет долю редкого класса в обеих выборках.

Числовые признаки стандартизованы с помощью \texttt{StandardScaler}, категориальные преобразованы в бинарные с помощью \texttt{OneHotEncoder}. Преобразования включены в конвейер каждой модели и обучаются только на тренировочной выборке. Балансировка классов не применялась: тестовая выборка сохраняет исходное распределение, а качество оценивается не только по доле верных ответов.

Обучены четыре классификатора: k-ближайших соседей с $k=5$, дерево решений, логистическая регрессия и случайный лес со 100 деревьями. Для воспроизводимости задано фиксированное зерно генератора. Подбор гиперпараметров и дополнительные задания не выполнялись.

\section{Оценка качества}

Для каждой модели рассчитаны Accuracy, Precision, Recall, Balanced accuracy и F1. Для положительного класса полнота показывает, какую долю пациентов класса 1 модель нашла, а Precision — какая доля предсказаний класса 1 верна. ROC AUC вычислена по вероятностным оценкам моделей. Результаты представлены в~\tabref{metrics}.

\suaitable[metrics]{Метрики классификации на тестовой выборке}
Модель | Accuracy | Precision | Recall | Balanced accuracy | F1 | ROC AUC
k-ближайших соседей | 0,987 | 1,000 | 0,200 | 0,600 | 0,333 | 0,792
Дерево решений | 0,968 | 0,154 | 0,200 | 0,591 | 0,174 | 0,591
Логистическая регрессия | 0,985 | 0,667 | 0,200 | 0,599 | 0,308 | 0,985
Случайный лес | 0,987 | 1,000 | 0,200 | 0,600 | 0,333 | 0,984

Матрицы неточностей приведены на~\figref{confusion_matrices}, ROC-кривые четырёх моделей — на~\figref{roc_curves}.

\suaiimg[0.86\textwidth]{confusion_matrices}{Матрицы неточностей классификаторов}

\suaiimg[0.75\textwidth]{roc_curves}{ROC-кривые классификаторов}

Точность Accuracy всех моделей высока, но сама по себе она не отражает качество распознавания редкого класса: если всегда предсказывать класс 0, доля верных ответов составила бы 98,4~\%. Каждая модель нашла только 2 из 10 положительных объектов тестовой выборки, поэтому Recall равен 0,2. k-ближайших соседей и случайный лес не дали ложноположительных предсказаний; у логистической регрессии одно такое предсказание, у дерева решений — 11.

По ROC AUC лучше логистическая регрессия (0,985) и случайный лес (0,984), но при стандартном пороге классификации полнота у них остаётся низкой. Высокое ROC AUC означает, что вероятностные оценки хорошо ранжируют объекты при разных порогах, но не означает, что выбранный порог уже подходит для поиска заболевших. В тесте всего 10 объектов класса 1, поэтому оценки чувствительны к ошибке даже на одном пациенте. Результаты показывают, что на этом разбиении ни одна модель не обеспечивает достаточную полноту положительного класса.

\suaiconclusion

В ходе работы подготовлен набор данных варианта 1 и обучены четыре требуемых классификатора. Категориальные и числовые признаки обработаны в конвейерах, а тестовые данные не участвовали в настройке преобразований. Из-за сильного дисбаланса классов Accuracy оказалась недостаточной метрикой для сравнения: при высокой точности все модели нашли лишь два из десяти положительных объектов тестовой выборки. Наибольший ROC AUC получен у логистической регрессии и случайного леса, однако низкая полнота не позволяет считать результат достаточным для диагностической задачи. Модели применимы только для учебного сравнения и не предназначены для медицинского использования.

\end{document}
